%% ====================================================================
\section{Hipótesis de partida}
\label{sec:hipotesis}
%% ====================================================================

\textbf{Pregunta de investigación.}
?`Pueden las dinámicas de Smith con una función de payoff multiplicativa y cota
inferior sigmoide de cardinalidad garantizar, de forma distribuida y sin coordinador
central, que una flota de AGV forme coaliciones de tamaño mínimo requerido para
el transporte cooperativo de cargas, y cómo cambia este comportamiento al variar
el rango de comunicación?

\textbf{Hipótesis~1.}
Bajo las dinámicas de Smith con la función de payoff propuesta, el sistema
converge a un equilibrio Nash en el que cada tarea $k$ tiene exactamente $n_k$ robots
asignados, en el caso simétrico balanceado ($N = Kn$, $n_k = n$, $\Psi \equiv 1$,
$R = \infty$). La condición suficiente de estabilidad local exige que la pendiente
$\beta$ de la función sigmoide pertenezca al intervalo $(\beta_{\min}, \beta_{\max})$.

\textbf{Hipótesis~2.}
La tasa de revisión espacialmente modulada $\mu_i(t)$ reduce significativamente
el número de cambios de asignación por robot durante el estado estacionario
respecto a una tasa constante $\mu_i = \mu_0$, sin degradar de forma relevante
el tiempo de convergencia de la coalición.

\textbf{Hipótesis~3.}
El método propuesto supera en tasa de factibilidad de formación de coaliciones
a la heurística voraz local y al consenso puro en todos los escenarios evaluados,
y la ventaja se mantiene cuando el rango de comunicación $R$ se reduce,
aunque la brecha disminuye con la conectividad.

\textbf{Hipótesis~4.}
La brecha de descentralización $\Delta J = (\Theta_{\text{cent}} -
\Theta_{\text{dist}}) / \max(\Theta_{\text{cent}}, \delta_0)$, con $\delta_0 = 10^{-3}$,
permanece acotada en los escenarios homogéneos y heterogéneos, con valores
aceptables para un entorno de almacén con comunicación imperfecta.

\subsection{Criterios cuantitativos de aceptación}
\label{subsec:criterios}

Los umbrales de aceptación que operacionalizan las hipótesis son los siguientes:

\begin{description}
  \item[H1 (convergencia Nash):] El número de episodios con convergencia
    ($|x_k(T_{\text{conv}}) - x_k^*| < 0{,}01$ para todo $k$) es $\geq 27$ de 30
    repeticiones en el Escenario~A balanceado ($N = Kn = 6$, $n_k = 2$).

  \item[H2 (revisión modulada):] $\Delta_{\text{chg}}$ del tratamiento T5 con
    tasa espacial es al menos un 40\,\% menor que con tasa constante $\mu_i = \mu_0$,
    sin que $T_{\text{coal}}$ aumente más de un 15\,\%. El 40\,\% se justifica
    porque, al valor nominal $\sigma_{\text{rev}} = 2\,\text{m}$, un robot a
    distancia $d = \sigma_{\text{rev}}$ tiene $\mu_i \approx 0{,}39\,\mu_0$; en
    estado estacionario la mayoría de robots están dentro de esa distancia,
    por lo que una reducción agregada del 40\,\% es conservadora.

  \item[H3 (factibilidad de coalición):] La tasa de factibilidad $F \geq 0{,}85$
    para T5 en todos los escenarios evaluados, incluyendo $R = 3\,\text{m}$.
    El umbral del 85\,\% refleja que una tasa de fallo del 15\,\% es recuperable
    mediante reintento en entornos reales; umbrales más laxos no serían informativos.

  \item[H4 (brecha de descentralización):] $\Delta J < 0{,}25$ en escenarios
    homogéneos y $\Delta J < 0{,}40$ en heterogéneos. Una brecha del 25\,\%
    significa que el sistema distribuido completa al menos 3 tareas de cada 4 que
    completaría un planificador centralizado con estado global perfecto.
\end{description}

Las hipótesis se contrastan mediante campaña de simulación con $m \geq 30$
repeticiones por tratamiento-escenario, bajo las mismas semillas, trayectorias
y parámetros iniciales para todos los métodos. La Hipótesis~1 se verifica
numéricamente en el escenario~A antes de ejecutar la campaña completa.
